\chapter{Modo explorador}
\label{sec:man-explorador}

\section{Qué es}

El modo explorador es la vista nodo-céntrica. Donde el modo caza
responde ``dada esta forma, ¿qué paths la tienen?'', el modo
explorador responde ``dada esta entidad, ¿a qué está conectada y
cuándo?''.

Se entra haciendo doble click sobre cualquier nodo de un resultado
de hunt, o desde \emph{Buscar}.

\section{Buscar entidades}

En la barra superior: escribir parte del valor (IP, hostname,
nombre de usuario) y filtrar opcionalmente por tipo. Resultados
aparecen con el conteo de eventos donde participa cada entidad.
Click $\to$ entra al modo explorador sobre ese nodo.

\section{Panel de detalle}

Para un nodo seleccionado:

\begin{itemize}
    \item \textbf{Propiedades}: tipo, valor, hash si aplica,
    timestamps de primera y última aparición.
    \item \textbf{Grados}: cuántas aristas entrantes y salientes.
    \item \textbf{Scores precomputados}: entity rarity,
    neighborhood concentration, temporal novelty (individuales del
    nodo, antes de ponderar).
    \item \textbf{Tags y notas}: lo que le pusiste.
    \item \textbf{Timeline}: histograma de actividad por hora /
    día.
\end{itemize}

\section{Expandir vecinos}

\emph{Expand}: muestra todos los nodos conectados, con la relación
y el timestamp del evento que produjo la arista. Para nodos con
muchos vecinos:

\emph{Expand grouped}: agrupa los vecinos por tipo de relación.
Primero ves las relaciones (\code{Connect} $\times$ 1523,
\code{DNS} $\times$ 87), después desplegás la que interese. Útil
cuando el nodo es un hub.

\section{Eventos del nodo}

Pestaña \emph{Events}: la lista cronológica de eventos donde
participa el nodo (como actor o target). Paginada. Cada entrada
tiene timestamp, los otros participantes, y atributos del evento.

Para nodos con miles de eventos (hub), la paginación cursor-based
evita OOM. Podés filtrar por ventana temporal o por tipo de
relación.

\section{Subgrafo}

Selección múltiple: ctrl+click sobre varios nodos, después
\emph{Subgraph}. Abre una vista del subgrafo inducido — todas las
aristas entre esos nodos. Exportable a GraphSON o JSONL para
compartir o diagramar en herramientas externas.

\section{Cuándo pasar de caza a exploración}

\begin{itemize}
    \item Encontraste un path con score alto y querés entender
    quién más participa alrededor de esas entidades.
    \item Un path sospechoso tiene un nodo hub y querés ver de dónde
    viene la concentración.
    \item Querés reconstruir narrativa temporal: qué pasó primero,
    qué vino después.
\end{itemize}

\section{Cuándo volver a caza}

\begin{itemize}
    \item Al explorar una IP externa sospechosa, ves que hay otros
    3 hosts que también se conectaron allí dentro de 2 horas —
    escribís un hunt nuevo con esa forma.
    \item El patrón que descubriste manualmente se puede generalizar
    a hipótesis reutilizable.
\end{itemize}

La iteración caza $\to$ exploración $\to$ caza es el modo normal
de trabajo.
